% Generated draft: requires mathematical and source review.
\documentclass[11pt]{article}
\usepackage[margin=25mm]{geometry}
\usepackage{fontspec}
\setmainfont{Latin Modern Roman}
\usepackage{amsmath,amssymb,mathtools,unicode-math}
\setmathfont{Latin Modern Math}
\usepackage{enumitem,needspace,xurl,hyperref}
\setlength{\parindent}{0pt}
\setlength{\parskip}{5pt}
\setlength{\emergencystretch}{3em}
\newcommand{\N}{\mathbb{N}}
\newcommand{\Z}{\mathbb{Z}}
\newcommand{\Q}{\mathbb{Q}}
\newcommand{\R}{\mathbb{R}}
\newcommand{\C}{\mathbb{C}}
\newcommand{\F}{\mathbb{F}}
\newcommand{\A}{\mathbb{A}}
\newcommand{\PP}{\mathbb{P}}
\newcommand{\E}{\mathbb{E}}
\newcommand{\eps}{\varepsilon}
\newcommand{\dd}{\,\mathrm{d}}
\begin{document}
\section*{UnsolvedMath ID: 30005032}
\textbf{Source problem code:} catalogue entry
\subsection{Definitions and mathematical statement}
Let \(W\) denote the symmetric stable driving process referred to in the source, and let \(b\) denote its distributional drift. A distributional drift is interpreted as a distribution in the spatial variable, rather than necessarily a function that can be evaluated along a sample path. The source concerns a linear equation labelled \((L)\) and a McKean--Vlasov equation labelled \((NL)\). In the latter setting, the drift depends on the distribution of the solution; write \(\mu_t=\mathcal L(X_t)\) for this distribution under the usual McKean--Vlasov interpretation.

For orientation only, their intended structure can be expressed schematically as
\[
X_t=\xi+A_t+W_t,
\]
where \(\xi\) is the initial random variable and \(A\) is a drift contribution that must be defined by the sought solution concept. For an ordinary function-valued drift in the linear case, the corresponding formal expression would be
\[
A_t=\int_0^t b(r,X_r)\,\dd r.
\]
For a distributional drift, this expression is not assumed to be defined. In the nonlinear case, the corrupted extract suggests an interaction involving \(b\) and \(\mu_r\), but does not reliably specify that interaction. These displays describe the intended structure and are not a recovered transcription of \((L)\) and \((NL)\).

\textbf{Problem.} Specify a notion of solution and a corresponding interpretation of the drift contribution for the source's linear and McKean--Vlasov equations with symmetric stable noise and distributional drift. Determine hypotheses under which this interpretation yields existence and uniqueness of the dynamics, including regimes below the relevant Young threshold, and state precisely the sense of uniqueness obtained.

\textbf{Unresolved qualifications.} The supplied record does not specify the spatial dimension, stability index or spectral characteristics of the noise, time interval, admissible initial laws, drift regularity spaces, time regularity, nonlinear interaction, or the meaning and numerical value of the Young threshold. It also does not specify whether uniqueness means uniqueness in law, pathwise uniqueness, or another property of the dynamics. The original extract omits \((L)\) and corrupts \((NL)\). Consequently, no fully quantified well-posedness assertion for a definite class of equations can be recovered from this record alone; in particular, the question must not be read as asserting well-posedness for every distributional drift.
\subsection{Short English statement}
Find a rigorous way to define and uniquely determine the linear and McKean--Vlasov dynamics driven by symmetric stable noise when the drift is a distribution, including below the applicable Young threshold. The supplied equations and hypotheses are incomplete, so the precise admissible class and meaning of uniqueness require clarification.
\subsection{Sources}
\begin{itemize}
\item Supplied local UnsolvedMath registry record: permanent ID \(30005032\), title \emph{Solution Concepts for Stable SDEs with Distributional Drift}, source problem code \texttt{catalogue entry}, in the Oberwolfach Reports Open Problems collection. No individual proposer is identified; the supplied proposed year is \(2022\).
\item Supplied report citation: \emph{Mini-Workshop: Regularization by Noise: Theoretical Foundations, Numerical Methods and Applications} (2023), Oberwolfach Reports, pp.\ 457--488, DOI \texttt{10.4171/owr/2022/9}. \url{https://doi.org/10.4171/owr/2022/9}. The background describes these reports as records of workshops and their open-problem discussions.
\item Supplied literature source labelled ``Stable driven McKean Vlasov equations'': \url{https://arxiv.org/abs/2205.11866}. The record supplies no authors, theorem references, or detailed hypotheses for this reference.
\item Literature assessment checked on \texttt{2026-08-22}: the supplied curation labels the problem open and states that martingale, virtual, Dirichlet, and related solution notions give well-posedness in restricted regularity regimes, while the general distributional-drift dynamics in the extract remain unresolved. This is supplied dated curation, not an independently verified assessment; its time-sensitive status and cited claims require independent verification.
\item The registry labels the statement \texttt{corrected\_verified} and reports that the official report was checked directly and that the cleaned statement restores definitions and removes extraction artifacts. These are supplied provenance claims, not verification performed for this draft; the supplied record still lacks the complete equations and quantitative hypotheses.
\end{itemize}
\end{document}
